\documentclass[11pt]{article}
% The preamble every project's paper shares: % The preamble every project's paper shares: % The preamble every project's paper shares: \input{../../preamble} from
% projects/<project>/paper/. Boards come from chessboard and xskak, plots from
% pgfplots; every number a paper prints comes from its generated.tex.
\usepackage[LSBC4,T1]{fontenc}
\usepackage[utf8]{inputenc}
\usepackage{lmodern}
\usepackage[margin=1in]{geometry}
\usepackage{microtype}
\usepackage{amsmath}
\usepackage{booktabs}
\usepackage{tabularx}
\usepackage{array}
\usepackage{xcolor}
\usepackage[obeyspaces,spaces]{url}
\usepackage{xskak}
\usepackage{pgfplots}
\pgfplotsset{compat=1.18}
\usepackage{caption}
\usepackage{subcaption}
\usepackage{enumitem}
\usepackage[section]{placeins}
\usepackage[hidelinks]{hyperref}

\setchessboard{boardfontencoding=LSBC4, boardfontfamily=skaknew, blackfieldcolor=black!22,
  whitefieldcolor=white, setfontcolors, showmover=false, label=true, labelfontsize=6pt,
  boardfontsize=14pt}

% A move or line in UCI, and a position in FEN, both breakable at a slash.
\newcommand{\uci}[1]{\texttt{#1}}
\newcommand{\fen}[1]{{\footnotesize\path{#1}}}
\newcolumntype{L}{>{\raggedright\arraybackslash}X}

\setlength{\parskip}{0.4em}
\setlength{\parindent}{0pt}
 from
% projects/<project>/paper/. Boards come from chessboard and xskak, plots from
% pgfplots; every number a paper prints comes from its generated.tex.
\usepackage[LSBC4,T1]{fontenc}
\usepackage[utf8]{inputenc}
\usepackage{lmodern}
\usepackage[margin=1in]{geometry}
\usepackage{microtype}
\usepackage{amsmath}
\usepackage{booktabs}
\usepackage{tabularx}
\usepackage{array}
\usepackage{xcolor}
\usepackage[obeyspaces,spaces]{url}
\usepackage{xskak}
\usepackage{pgfplots}
\pgfplotsset{compat=1.18}
\usepackage{caption}
\usepackage{subcaption}
\usepackage{enumitem}
\usepackage[section]{placeins}
\usepackage[hidelinks]{hyperref}

\setchessboard{boardfontencoding=LSBC4, boardfontfamily=skaknew, blackfieldcolor=black!22,
  whitefieldcolor=white, setfontcolors, showmover=false, label=true, labelfontsize=6pt,
  boardfontsize=14pt}

% A move or line in UCI, and a position in FEN, both breakable at a slash.
\newcommand{\uci}[1]{\texttt{#1}}
\newcommand{\fen}[1]{{\footnotesize\path{#1}}}
\newcolumntype{L}{>{\raggedright\arraybackslash}X}

\setlength{\parskip}{0.4em}
\setlength{\parindent}{0pt}
 from
% projects/<project>/paper/. Boards come from chessboard and xskak, plots from
% pgfplots; every number a paper prints comes from its generated.tex.
\usepackage[LSBC4,T1]{fontenc}
\usepackage[utf8]{inputenc}
\usepackage{lmodern}
\usepackage[margin=1in]{geometry}
\usepackage{microtype}
\usepackage{amsmath}
\usepackage{booktabs}
\usepackage{tabularx}
\usepackage{array}
\usepackage{xcolor}
\usepackage[obeyspaces,spaces]{url}
\usepackage{xskak}
\usepackage{pgfplots}
\pgfplotsset{compat=1.18}
\usepackage{caption}
\usepackage{subcaption}
\usepackage{enumitem}
\usepackage[section]{placeins}
\usepackage[hidelinks]{hyperref}

\setchessboard{boardfontencoding=LSBC4, boardfontfamily=skaknew, blackfieldcolor=black!22,
  whitefieldcolor=white, setfontcolors, showmover=false, label=true, labelfontsize=6pt,
  boardfontsize=14pt}

% A move or line in UCI, and a position in FEN, both breakable at a slash.
\newcommand{\uci}[1]{\texttt{#1}}
\newcommand{\fen}[1]{{\footnotesize\path{#1}}}
\newcolumntype{L}{>{\raggedright\arraybackslash}X}

\setlength{\parskip}{0.4em}
\setlength{\parindent}{0pt}

% Written by projects/judge-in-context/paper/generate.mjs; do not edit by hand.

\newcommand{\ThBlunder}{20}
\newcommand{\ThMistake}{10}
\newcommand{\ThInaccuracy}{5}
\newcommand{\ThGood}{1}
\newcommand{\ThKeptWon}{88}
\newcommand{\ThKeptWinning}{75}
\newcommand{\ShippedSlope}{0.00368208}
\newcommand{\LowSlope}{0.00245}
\newcommand{\RatingFloor}{1000}
\newcommand{\RatingSpan}{1300}
\newcommand{\RatingTop}{2300}
\newcommand{\Games}{60}
\newcommand{\Records}{3917}
\newcommand{\ExactPlies}{3880}
\newcommand{\NotExactPlies}{37}
\newcommand{\BookPlies}{254}
\newcommand{\ForcedPlies}{48}
\newcommand{\GradedPlies}{3578}
\newcommand{\GradedPlayer}{1783}
\newcommand{\GradedOpponent}{1795}
\newcommand{\Flips}{492}
\newcommand{\FlipShare}{13.8}
\newcommand{\FlipsPlayer}{250}
\newcommand{\FlipsOpponent}{242}
\newcommand{\FlipSharePlayer}{14.0}
\newcommand{\FlipShareOpponent}{13.5}
\newcommand{\MateSlipChanged}{9}
\newcommand{\PlayerRatingMin}{948}
\newcommand{\PlayerRatingMax}{1167}
\newcommand{\OpponentRatingMin}{513}
\newcommand{\OpponentRatingMax}{1178}
\newcommand{\CrossoverPawns}{5.1}
\newcommand{\FloorOld}{5.4}
\newcommand{\FloorNew}{8.1}
\newcommand{\CapOld}{3.0}
\newcommand{\CapNew}{4.5}
\newcommand{\AssertMaxDiff}{1.4\times 10^{-14}}
\newcommand{\EvidenceRecorded}{review-evidence/4}
\newcommand{\FlipOneIn}{seven}
\newcommand{\MainTcGames}{59}
\newcommand{\MainTcMinutes}{10}
\newcommand{\JudgeStudent}{an adult club player rated about 950--1150 on chess.com who plays 10-minute games}
\newcommand{\MatrixPlayer}{%
\begin{tabular}{l rrrrr r}
\toprule
shipped $\downarrow$ proposed $\rightarrow$ & blunder & mistake & inaccuracy & ok & good & total \\
\midrule
blunder & (37) & \textbf{29} & 0 & 0 & 0 & 66 \\
mistake & \textbf{4} & (50) & \textbf{48} & 0 & 0 & 102 \\
inaccuracy & \textbf{1} & \textbf{9} & (76) & \textbf{54} & 0 & 140 \\
ok & 0 & 0 & \textbf{16} & (439) & \textbf{78} & 533 \\
good & 0 & 0 & 0 & \textbf{11} & (931) & 942 \\
\midrule
total & 42 & 88 & 140 & 504 & 1009 & 1783 \\
\bottomrule
\end{tabular}
}
\newcommand{\MatrixOpponent}{%
\begin{tabular}{l rrrrr r}
\toprule
shipped $\downarrow$ proposed $\rightarrow$ & blunder & mistake & inaccuracy & ok & good & total \\
\midrule
blunder & (52) & \textbf{28} & 0 & 0 & 0 & 80 \\
mistake & \textbf{2} & (45) & \textbf{50} & 0 & 0 & 97 \\
inaccuracy & 0 & \textbf{10} & (75) & \textbf{47} & 0 & 132 \\
ok & 0 & \textbf{2} & \textbf{12} & (500) & \textbf{76} & 590 \\
good & 0 & 0 & 0 & \textbf{15} & (881) & 896 \\
\midrule
total & 54 & 85 & 137 & 562 & 957 & 1795 \\
\bottomrule
\end{tabular}
}
\newcommand{\JudgeMoves}{25}
\newcommand{\RunsAgree}{23}
\newcommand{\ShippedRunOne}{13}
\newcommand{\ShippedRunTwo}{15}
\newcommand{\ProposedRunOne}{12}
\newcommand{\ProposedRunTwo}{10}
\newcommand{\UnsureRunOne}{0}
\newcommand{\UnsureRunTwo}{0}
\newcommand{\JudgeTable}{%
\begin{tabular}{lrr}
\toprule
 & run 1 & run 2 \\
\midrule
shipped verdict preferred & 13 & 15 \\
proposed verdict preferred & 12 & 10 \\
unsure & 0 & 0 \\
\bottomrule
\end{tabular}
}
\newcommand{\KindTable}{%
\begin{tabular}{lrrrr}
\toprule
shown label: shipped $\to$ proposed & moves & both runs: shipped & both runs: proposed & split \\
\midrule
Inaccuracy $\to$ Good & 7 & 2 & 5 & 0 \\
Mistake $\to$ Inaccuracy & 7 & 4 & 1 & 2 \\
Blunder $\to$ Mistake & 4 & 2 & 2 & 0 \\
Good $\to$ Inaccuracy & 2 & 2 & 0 & 0 \\
Inaccuracy $\to$ Mistake & 2 & 1 & 1 & 0 \\
Mistake $\to$ Blunder & 2 & 1 & 1 & 0 \\
Inaccuracy $\to$ Blunder & 1 & 1 & 0 & 0 \\
\bottomrule
\end{tabular}
}
\newcommand{\ExampleCount}{4}
\newcommand{\ExampleACaption}{Inaccuracy on the shipped curve, Good on the proposed one: the judge preferred the proposed label in both runs.}
\newcommand{\ExampleA}{% Written by projects/judge-in-context/paper/generate.mjs; do not edit by hand.
\begin{minipage}[t]{0.4\textwidth}\vspace{0pt}
\chessboard[setfen={rnbq1rk1/ppp1bppp/3p1p2/8/3P4/1QP1P3/PP3PPP/RN2KBNR b KQ - 0 6}, boardfontsize=13pt, pgfstyle=straightmove, arrow=to, linewidth=0.07em, color=red!75!black, markmoves={b8-c6}, pgfstyle=straightmove, arrow=to, linewidth=0.07em, color=green!50!black, markmoves={d6-d5}]
\end{minipage}\hfill
\begin{minipage}[t]{0.56\textwidth}\vspace{0pt}\small
\textbf{Played:} 6\ldots{}\symknight{}c6, Black to move; ratings 1076 (the player) and 1074.\par
\textbf{Engine:} +0.29 with the best move, $-$0.33 after yours.\par
\textbf{Best line:} 6\ldots{}d5 7.~\symknight{}d2 a5 8.~h4 a4 9.~\symqueen{}c2\par\textbf{UCI:} \uci{b8c6}; best line \uci{d6d5 b1d2 a7a5 h2h4 a5a4 b3c2}\par
\textbf{Shipped:} Inaccuracy. \textbf{Proposed:} Good. \textbf{Judge:} the proposed verdict, in both runs.\par
\textbf{Reason (run 1):} ``Nc6 is a normal developing move and the position stays balanced; the engine's small preference for ...d5 does not matter at this level.''\par
\textbf{FEN:} \fen{rnbq1rk1/ppp1bppp/3p1p2/8/3P4/1QP1P3/PP3PPP/RN2KBNR b KQ - 0 6}
\end{minipage}
}
\newcommand{\ExampleBCaption}{Mistake on the shipped curve, Inaccuracy on the proposed one: the judge preferred the shipped label in both runs.}
\newcommand{\ExampleB}{% Written by projects/judge-in-context/paper/generate.mjs; do not edit by hand.
\begin{minipage}[t]{0.4\textwidth}\vspace{0pt}
\chessboard[setfen={r2qkb1r/ppp2ppp/2n2n2/3ppb2/2PP1B2/4PN2/PP3PPP/RN1QKB1R w KQkq - 0 6}, boardfontsize=13pt, pgfstyle=straightmove, arrow=to, linewidth=0.07em, color=red!75!black, markmoves={c4-d5}, pgfstyle=straightmove, arrow=to, linewidth=0.07em, color=green!50!black, markmoves={f3-e5}]
\end{minipage}\hfill
\begin{minipage}[t]{0.56\textwidth}\vspace{0pt}\small
\textbf{Played:} 6.~cxd5, White to move; ratings 1040 (the player) and 1052.\par
\textbf{Engine:} +1.05 with the best move, $-$0.35 after yours.\par
\textbf{Best line:} 6.~\symknight{}xe5 \symknight{}b4 7.~\symknight{}a3 c6 8.~\symbishop{}e2 \symbishop{}d6\par\textbf{UCI:} \uci{c4d5}; best line \uci{f3e5 c6b4 b1a3 c7c6 f1e2 f8d6}\par
\textbf{Shipped:} Mistake. \textbf{Proposed:} Inaccuracy. \textbf{Judge:} the shipped verdict, in both runs.\par
\textbf{Reason (run 1):} ``Black's e5 pawn was simply hanging to Nxe5, and cxd5 instead allows ...exf4/...Bb4+ play where White is no longer better, so a full pawn's worth is thrown away.''\par
\textbf{FEN:} \fen{r2qkb1r/ppp2ppp/2n2n2/3ppb2/2PP1B2/4PN2/PP3PPP/RN1QKB1R w KQkq - 0 6}
\end{minipage}
}
\newcommand{\ExampleCCaption}{Blunder on the shipped curve, Mistake on the proposed one: the judge preferred the shipped label in both runs.}
\newcommand{\ExampleC}{% Written by projects/judge-in-context/paper/generate.mjs; do not edit by hand.
\begin{minipage}[t]{0.4\textwidth}\vspace{0pt}
\chessboard[setfen={r1bqkbnr/1p3ppp/p1n1p3/1Bpp4/3P1B2/4PN2/PPP2PPP/RN1QK2R w KQkq - 0 6}, boardfontsize=13pt, pgfstyle=straightmove, arrow=to, linewidth=0.07em, color=red!75!black, markmoves={b5-a4}, pgfstyle=straightmove, arrow=to, linewidth=0.07em, color=green!50!black, markmoves={b5-c6}]
\end{minipage}\hfill
\begin{minipage}[t]{0.56\textwidth}\vspace{0pt}\small
\textbf{Played:} 6.~\symbishop{}a4, White to move; ratings 1102 (the player) and 1105.\par
\textbf{Engine:} +0.31 with the best move, $-$2.51 after yours.\par
\textbf{Best line:} 6.~\symbishop{}xc6+ bxc6 7.~\symknight{}c3 cxd4 8.~exd4 f6\par\textbf{UCI:} \uci{b5a4}; best line \uci{b5c6 b7c6 b1c3 c5d4 e3d4 f7f6}\par
\textbf{Shipped:} Blunder. \textbf{Proposed:} Mistake. \textbf{Judge:} the shipped verdict, in both runs.\par
\textbf{Reason (run 1):} ``Ba4 walks into ...b5 and ...c4, which traps the bishop and costs a piece for a pawn: decisive material in the opening.''\par
\textbf{FEN:} \fen{r1bqkbnr/1p3ppp/p1n1p3/1Bpp4/3P1B2/4PN2/PPP2PPP/RN1QK2R w KQkq - 0 6}
\end{minipage}
}
\newcommand{\ExampleDCaption}{Good on the shipped curve, Inaccuracy on the proposed one: the judge preferred the shipped label in both runs.}
\newcommand{\ExampleD}{% Written by projects/judge-in-context/paper/generate.mjs; do not edit by hand.
\begin{minipage}[t]{0.4\textwidth}\vspace{0pt}
\chessboard[setfen={8/p3r1k1/1p6/5K1p/P3p2P/1P2R1P1/2r5/8 b - - 1 45}, boardfontsize=13pt, pgfstyle=straightmove, arrow=to, linewidth=0.07em, color=red!75!black, markmoves={c2-c5}, pgfstyle=straightmove, arrow=to, linewidth=0.07em, color=green!50!black, markmoves={c2-f2}]
\end{minipage}\hfill
\begin{minipage}[t]{0.56\textwidth}\vspace{0pt}\small
\textbf{Played:} 45\ldots{}\symrook{}c5+, Black to move; ratings 1105 (the player) and 1099.\par
\textbf{Engine:} you mate in 3 with the best move, +5.91 after yours.\par
\textbf{Best line:} 45\ldots{}\symrook{}f2+ 46.~\symrook{}f3 \symrook{}xf3+ 47.~\symking{}g5 \symrook{}e5\#\par\textbf{UCI:} \uci{c2c5}; best line \uci{c2f2 e3f3 f2f3 f5g5 e7e5}\par
\textbf{Shipped:} Good. \textbf{Proposed:} Inaccuracy. \textbf{Judge:} the shipped verdict, in both runs.\par
\textbf{Reason (run 1):} ``Rf2+ mated in three, but after Rc5+ Black is still a full rook up and easily winning, so nothing important was given away.''\par
\textbf{FEN:} \fen{8/p3r1k1/1p6/5K1p/P3p2P/1P2R1P1/2r5/8 b - - 1 45}
\end{minipage}
}
\newcommand{\EvidenceRerun}{review-evidence/5}
\newcommand{\RerunFlips}{492}
\newcommand{\RerunFlipShare}{13.8}
\newcommand{\RerunFlipsPlayer}{247}
\newcommand{\RerunFlipsOpponent}{245}
\newcommand{\ReproTable}{%
\begin{tabular}{lrr}
\toprule
 & recorded & rerun in the lab \\
\midrule
engine evidence & review-evidence/4 & review-evidence/5 \\
games & 60 & 60 \\
plies with exact scores & 3880 & 3880 \\
graded plies & 3578 & 3578 \\
flips & 492 (13.8\%) & 492 (13.8\%) \\
the player's flips & 250 of 1783 & 247 of 1783 \\
the opponents' flips & 242 of 1795 & 245 of 1795 \\
the copy of the judge matches the app's & yes & yes \\
\bottomrule
\end{tabular}
}
\newcommand{\AppCommit}{\texttt{d8b7712}}
\newcommand{\RecordsCommit}{\texttt{b3b827d}}


\title{Should a move's verdict follow the mover's rating?\\[0.3em]
\large A rating-aware win-probability curve, put to a blind judge}
\author{Chess Coach research lab}
\date{}

\begin{document}
\maketitle

\begin{abstract}
The coaching app labels each move (good, ok, inaccuracy, mistake, blunder) from how much
of the mover's winning chances it gives away, read on the win-probability curve Lichess
uses. A natural objection is that the curve is calibrated on stronger players: near a
level position, a pawn decides fewer games between lower-rated players. We tested a
proposal that flattens the curve for lower-rated movers, with a slope of \LowSlope{} per
centipawn at or below \RatingFloor{} rising linearly to the shipped \ShippedSlope{} at
\RatingTop{}. Over \Games{} of the player's games, the proposal changes \Flips{} of the
\GradedPlies{} graded verdicts (\FlipShare\%). A judge model, shown \JudgeMoves{} of the
player's own flipped moves with the two labels in random order and unnamed, preferred the
shipped label \ShippedRunOne{} and \ShippedRunTwo{} times in two runs, and the proposed one
\ProposedRunOne{} and \ProposedRunTwo{} times. The proposal moves one verdict in \FlipOneIn{}
without making verdicts fairer on balance, and was not adopted. The judge's reasons point
elsewhere: to what a move gives away (material, a mate) rather than to the size of the swing.
\end{abstract}

\section{The question}

A game review is only useful if its verdicts match what the player can learn from. The
app's judge reads the mover's win probability before and after the move and labels the
move by the drop. The curve that turns an engine evaluation into a win probability is the
one Lichess fits on its own games, and it does not depend on who is playing. For a player
at this level, a half-pawn concession in a level position rarely decides the game, so a
curve for that level might be flatter near zero, and verdicts there milder.

The question here is whether such a curve gives fairer verdicts \emph{to this player}: an
adult club player, rated \PlayerRatingMin--\PlayerRatingMax{} on chess.com across the \Games{} games
studied here, \MainTcGames{} of them played at \MainTcMinutes{} minutes a side.

\section{The judge and the proposal}

The win probability of the mover, for an evaluation of $x$ centipawns from the mover's side,
is
\[
  W_k(x) \;=\; 50 + 50\left(\frac{2}{1+e^{-kx}} - 1\right),
\]
with $k = \ShippedSlope$ in the app. The judge takes the drop
$\Delta = W_k(x_\text{before}) - W_k(x_\text{after})$ in percentage points and labels the
move a blunder at $\Delta \ge \ThBlunder$, a mistake at $\ge \ThMistake$, an inaccuracy at
$\ge \ThInaccuracy$, good at $\le \ThGood$, and ok otherwise. Two clamps protect a decided
game: if the mover still has at least \ThKeptWon\% after the move, any label but good becomes
ok, and at \ThKeptWinning\% or more, a blunder or mistake becomes an inaccuracy. The judge
reads ok and good alike as a good move, so a flip between them does not change the label
the player sees.

The proposal changes only the slope, by the mover's rating $r$:
\[
  k(r) \;=\; \LowSlope + (\ShippedSlope - \LowSlope)\cdot
  \operatorname{clamp}\!\left(\frac{r - \RatingFloor}{\RatingSpan},\, 0,\, 1\right).
\]
Thresholds and clamps stay as shipped. Figure~\ref{fig:curves} shows the two curves for a
mover at or below \RatingFloor.

\begin{figure}[!htbp]
\centering
\begin{tikzpicture}
\begin{axis}[width=0.8\textwidth, height=6cm, xlabel={evaluation, pawns (mover's side)},
  ylabel={mover's win \%}, xmin=-10, xmax=10, ymin=0, ymax=100, grid=major,
  legend pos=south east, legend cell align=left]
\addplot[thick, black, domain=-10:10, samples=201] {50 + 50*(2/(1+exp(-\ShippedSlope*100*x)) - 1)};
\addlegendentry{shipped, $k=\ShippedSlope$}
\addplot[thick, red!70!black, dashed, domain=-10:10, samples=201] {50 + 50*(2/(1+exp(-\LowSlope*100*x)) - 1)};
\addlegendentry{proposed at $r \le \RatingFloor$, $k=\LowSlope$}
\end{axis}
\end{tikzpicture}
\caption{The two win-probability curves. Near level the proposed curve counts each pawn as
fewer win points, so a drop there reads milder; further out it counts each pawn as more.}
\label{fig:curves}
\end{figure}

\paragraph{Where the curves part.} The two curves change win probability at the same rate
per pawn at about \CrossoverPawns{} pawns from level on either side. Nearer level, the
proposed curve moves less per pawn and its verdicts are milder; further out it moves more
and they are harsher. The clamps move too: the \ThKeptWon\% floor needs about
$+\FloorOld$ pawns on the shipped curve and $+\FloorNew$ on the proposed one, and the
\ThKeptWinning\% cap needs $+\CapOld$ and $+\CapNew$. A move that leaves the mover between
those levels is protected on the shipped curve and not on the proposed one.

\section{How many verdicts change}

The data are the phone-profile reviews of \Games{} of the player's games
(\Records{} moves, \ExactPlies{} with an exact engine score before and after; the other
\NotExactPlies{} end the game with mate, which the outcome grades). Book and forced moves
(\BookPlies{} and \ForcedPlies) are not graded, which leaves \GradedPlies{} graded moves,
\GradedPlayer{} by the player and \GradedOpponent{} by the opponents. Player ratings in the
set run from \PlayerRatingMin{} to \PlayerRatingMax{} and opponent ratings from
\OpponentRatingMin{} to \OpponentRatingMax, so nearly every mover sits at or near the
proposal's floor.

Every graded move is judged twice by a copy of the app's judge that takes the slope as a
parameter. On the shipped slope the copy reproduces the app's own judge on all
\ExactPlies{} exact-score moves, with the largest win-probability difference
$\AssertMaxDiff$, so the comparison isolates the curve.

The proposal changes \Flips{} verdicts (\FlipShare\% of graded moves): \FlipsPlayer{} of the
player's (\FlipSharePlayer\%) and \FlipsOpponent{} of the opponents' (\FlipShareOpponent\%).
Table~\ref{tab:matrix} shows where they go.

\begin{table}[!htbp]
\centering\small
\begin{subtable}{\textwidth}\centering
\MatrixPlayer
\caption{The player's moves.}
\end{subtable}

\vspace{0.8em}
\begin{subtable}{\textwidth}\centering
\MatrixOpponent
\caption{The opponents' moves.}
\end{subtable}
\caption{Verdicts on the shipped curve (rows) and on the proposed one (columns). Unchanged
verdicts are in parentheses; flips are in bold.}
\label{tab:matrix}
\end{table}

Most flips are milder near level (an inaccuracy that becomes ok, a mistake that becomes an
inaccuracy, a blunder that becomes a mistake). The harsher ones come from positions already
far from level, where the proposed curve moves more per pawn and its clamps protect less.

\section{What a judge preferred}

\paragraph{Protocol.} The \JudgeMoves{} moves are the player's own flipped moves whose shown
label changes (the moves on the app's Coach Calibration page). For each, a judge model was
given the position before the move, the move, the engine's evaluation with the best move and
after the move played, the engine's best move and line, both players' ratings, and the two
labels as X and Y in random order, unnamed. It was asked, as an experienced coach whose student
is ``\JudgeStudent'', which label matches how bad the move really was at this level, or to say it could not tell, with one sentence of reason. It answered twice,
independently, with the order of the labels reversed.

\paragraph{Result.} Table~\ref{tab:judge} shows the preferences. The two runs gave the same
answer on \RunsAgree{} of \JudgeMoves{} moves. Neither run prefers the proposed curve.

\begin{table}[!htbp]
\centering\small
\begin{subtable}{\textwidth}\centering
\JudgeTable
\caption{Each run.}
\end{subtable}

\vspace{0.8em}
\begin{subtable}{\textwidth}\centering
\KindTable
\caption{By the kind of change, as the player sees the labels.}
\end{subtable}
\caption{Which verdict the judge preferred.}
\label{tab:judge}
\end{table}

By kind, the judge took the proposed side where a small concession near level becomes a
good move, and kept the shipped side where the move simply hangs material, where a piece is
trapped, and where the game was already decided and the proposal would raise the verdict.
Figures~\ref{fig:exa}--\ref{fig:exd} show one move of each of the four commonest kinds, with
the move played (red arrow) and the engine's best move (green).

\begin{figure}[!htbp]\ExampleA\caption{\ExampleACaption}\label{fig:exa}\end{figure}
\begin{figure}[!htbp]\ExampleB\caption{\ExampleBCaption}\label{fig:exb}\end{figure}
\begin{figure}[!htbp]\ExampleC\caption{\ExampleCCaption}\label{fig:exc}\end{figure}
\begin{figure}[!htbp]\ExampleD\caption{\ExampleDCaption}\label{fig:exd}\end{figure}

\section{Decision, and what the reasons suggest}

The proposal would move one verdict in \FlipOneIn{}, and in neither run did the judge prefer the
proposed verdict more often than the shipped one, so it was not adopted. The reasons the judge gives are about consequences, not
the size of the swing: a move that hangs material near level is a mistake whatever the curve
says; a small positional concession near level is a normal move; and an error in a decided
position matters less, which the shipped clamps already express. Those are the theses a next
round would test, on a fresh sample, with the player's own answers as the key.

\section{Reproduction in the lab}

The lab recomputes every number of the judge's table from the judge's recorded answers and
each run's order key (\texttt{tally.mjs}), exactly. The flips are rerun on engine evidence
the lab rebuilt on the pinned app (\AppCommit) from the same games
(\texttt{reproduce.mjs}); the recorded run used evidence the app wrote before a later
change to its review pipeline (\EvidenceRecorded, against \EvidenceRerun{} now), so the
scores differ in detail. Table~\ref{tab:repro} sets the two side by side; the records are
read at \RecordsCommit.

\begin{table}[!htbp]
\centering\small
\ReproTable
\caption{The flips as recorded and as rerun in the lab.}
\label{tab:repro}
\end{table}

\section{Limitations}

\begin{itemize}[nosep]
\item The judge is a model, not the player. The player's own answers on the Calibration page
are the key this check is waiting for.
\item \JudgeMoves{} moves from one player's games is a small sample: it can show that the
proposal is not clearly fairer, not that it is clearly worse.
\item Only the slope was varied. The judge's reasons suggest that a fairer verdict needs
something the curve cannot see, such as whether the move gives away material.
\end{itemize}

\end{document}
